% In this file you should put all LaTeX macros and settings to be used both by
% the pdf version and the web version.
% This should be most of your macros.

% The theorem-like environments defined below are those that appear by default
% in the dependency graph. See the README of leanblueprint if you need help to
% customize this.
% The configuration below use the theorem counter for all those environments
% (this is what the [theorem] arguments mean) and never resets it.
% If you want for instance to number them within chapters then you can add
% [chapter] at the end of the next line.
\newtheorem{theorem}{Theorem}
\newtheorem{proposition}[theorem]{Proposition}
\newtheorem{lemma}[theorem]{Lemma}
\newtheorem{corollary}[theorem]{Corollary}

\theoremstyle{definition}
\newtheorem{definition}[theorem]{Definition}

% Standard fields
\newcommand{\C}{\mathbb{C}}
\newcommand{\Q}{\mathbb{Q}}
\newcommand{\Z}{\mathbb{Z}}
\newcommand{\N}{\mathbb{N}}
\newcommand{\PP}{\mathbb{P}}
\newcommand{\K}{\mathrm{K}}
\newcommand{\Sym}{\mathrm{Sym}}
\newcommand{\Pic}{\mathrm{Pic}}
\newcommand{\Gr}{\mathrm{Gr}}

% Calligraphic letters
\newcommand{\cO}{\mathcal{O}}
\newcommand{\cM}{\mathcal{M}}
\newcommand{\cS}{\mathcal{S}}
\newcommand{\cG}{\mathcal{G}}
\newcommand{\cZ}{\mathcal{Z}}
\newcommand{\cL}{\mathcal{L}}
\newcommand{\cC}{\mathcal{C}}
\newcommand{\oM}{\overline{\mathcal{M}}}

% The Picard stack
\newcommand{\Picabs}{\mathfrak{Pic}}

% Operational Chow
\newcommand{\CHop}{\mathsf{CH}_{\mathsf{op}}}

% Double ramification cycle
\newcommand{\DRop}{\mathsf{DR}^{\mathsf{op}}}

% Pixton's formula class
\newcommand{\Pix}{\mathsf{P}}

% Graph combinatorics
\newcommand{\GG}{\mathsf{G}}  % set of prestable graphs
\newcommand{\DGG}{\mathsf{DG}} % set of decorated prestable graphs
\newcommand{\Aut}{\mathrm{Aut}}
\newcommand{\val}{\mathrm{val}}

% Tautological classes
\newcommand{\psi}{\psi}
\newcommand{\xii}{\xi}
\newcommand{\etaab}{\eta}

% Chow groups
\newcommand{\Chow}{\mathsf{CH}}

% Moduli stacks
\newcommand{\fM}{\mathfrak{M}}
\newcommand{\fC}{\mathfrak{C}}

% Sets of vertices, half-edges, edges, legs
\newcommand{\V}{\mathrm{V}}
\newcommand{\HH}{\mathrm{H}}
\newcommand{\E}{\mathrm{E}}
\newcommand{\LL}{\mathrm{L}}

% Half-edge first Chern class
\newcommand{\field}{K}
